\documentclass[main.tex]{subfiles}

\begin{document}
\newpage
\section{Matrix Ordered Spaces}
\label{sec:matrix-ordered-spaces}

\begin{motivation}{\large Summary}{mot: matrix ordered spaces}
In this lecture, we will cover the following topics:
\begin{enumerate}
    \item The basic structure of matrix ordered $\star$-vector spaces and the
    completely positive maps between them.
    \item The Choi--Effros representation theorem, which realizes every
    abstract operator system as a concrete operator system.
\end{enumerate}
The principal source is Choi and Effros,
\emph{Injectivity and Operator Spaces}(\cite{choi-effros-injectivity-1977}), though we do make use of other books which discuss operator systems. In \Cref{subsec:matrix-ordered-spaces-I}, we include a brief introduction to (pre)-ordered vector spaces based partly on \cite{jameson1970ordered}. 
\end{motivation}

Here are the less common notational conventions used in this lecture.

\begin{filingbox}[NavyBlue]{Conventions}
If $V$ is a complex $\star$-vector space, then $M_n(V)$ denotes the
$\star$-vector space of $n$-by-$n$ matrices over $V$.  The symbols
$\mathcal S$ and $\mathcal T$ are reserved for operator systems, and all
unital completely positive maps are complex linear.
\end{filingbox}

\newpage

Let's start with some motivation.

\begin{motiv}
Self-testing statements are ultimately phrased in terms of positive elements,
states, and completely positive maps at every matrix level.  Matrix ordered
spaces isolate precisely this structure.  The abstract theory is useful only
if its objects can be represented concretely inside operator algebras; the
Choi--Effros theorem provides that bridge and is therefore the endpoint of
this lecture.
\end{motiv}

\end{document}
